\graphicspath{ {./images/5_Normalized_data} }
\newpage

\section{Normalized data}
Following analyte curation, total area normalization is automatically conducted on the data per glycosylation site. 
A table displaying the normalized data is generated, including the following information:

\begin{itemize}[topsep=0pt, itemsep=0pt]
	\item Sample names
	\item Sample types
	\item Sample IDs
	\item Plate well positions (if applicable).
	\item The number of replicates for each sample ID.
	\item Metadata (if applicable).
	\item Columns showing the total spectrum intensities for per glycosylation site.
	\item For each analyte, a column displaying the relative abundance in percentages.
	Analytes associated to the same glycosylation site will collectively sum to 100\%.
\end{itemize}

\begin{optionalbox}
	Above the table, there is an option to normalize charge states separately instead of combining them per analyte.
\end{optionalbox}

The normalized data is visually represented in one or more heatmaps, with glycan compositions on the $x$-axis and glycosylation sites on the $y$-axis by default.
Each color in the heatmap corresponds to the median relative abundance for a specific glycopeptide, calculated across all samples.
There is an option to exclude sample types from this calculation.
If analyte curation was performed per biological group, there will be an option to separate heatmaps per biological group (enabled by default).
\begin{center}
	\includegraphics[scale=0.42]{heatmap.png}
\end{center}

Alternatively, the individual samples can be plotted on the $y$-axis.
In this case, a heatmap is generated for each glycosylation site, with the ability to exclude sample types. 
If curation was performed per biological groups, the heatmaps can be generated per group.

